%% =====================================================================================
%%  FLOWCHART STYLE GUIDE  --  house style for one-page process flowcharts
%%  Rev. 1, 2026-10-14.  One file: instructions (Part 1), the drawing engine (Part 2),
%%  the content you write (Part 3). Compiles as-is with pdfLaTeX (MiKTeX or TeX Live).
%% =====================================================================================
%%
%%  PART 1 -- INSTRUCTIONS FOR THE AI THAT WRITES A FLOWCHART FROM THIS FILE
%%
%%  You are given this file and a description of a process. Produce ONE complete .tex
%%  file in ONE code block: it starts at \documentclass and ends at \end{document}.
%%  Do not include Part 1 (these comment lines above \documentclass). No text outside the
%%  code block except, if needed, one or two lines of notes after it.
%%
%%  THE ONE RULE: copy Part 2 (from \documentclass to the line "\begin{document}")
%%  exactly, character for character. Never shorten it, never "summarise" it, never add
%%  packages, \tikzset, \node, \draw, coordinates or lengths. The engine lays the drawing
%%  out by itself: it measures every box, spaces the columns, routes every arrow and
%%  fits the page. Your job is content only: Part 3, the lines between \begin{document}
%%  and \end{document}. Write them with the commands below and nothing else.
%%
%%  CONTENT COMMANDS (each on its own line; keep the braces balanced)
%%
%%  \ChartTitle{...}          Required. The drawing's title, 2 to 9 words, Title Case.
%%  \ChartSubtitle{...}       Optional. One italic line: what the chart shows, for whom.
%%  \Summary{...}{...}        Optional. Key facts as bullets in two columns (2 to 3 bullets
%%                            each, up to 10 words each). Second column may be empty: {}.
%%  \KeyBox{Heading}{...}     Optional. A framed box top right (e.g. Before you start,
%%                            Scope, Rules): 2 to 4 bullets, up to 14 words each.
%%  \Station{Heading}{...}    Required, 2 to 6 of them, in process order. One step of the
%%                            process: a green numbered header (2 to 4 words) and a white
%%                            box of 2 to 6 bullets (up to 12 words each). Numbering is
%%                            automatic. Stations are drawn left to right with arrows.
%%  \DocStation{Heading}{Lead line}{\doc{name}{tag}...}
%%                            A station whose box shows files or documents instead of
%%                            bullets (inputs at the start, outputs at the end): a short
%%                            lead line, then 1 to 4 \doc{file name}{tag} items, e.g.
%%                            \doc{Order Intake*.xlsx}{required}. Keep names under 24
%%                            characters; the tag may be empty: \doc{name}{}.
%%  \Gate{N}{Question?}{Yes}{No}{Outcome heading}{...}
%%                            Optional, at most one per two stations. A decision diamond
%%                            under station N (N must be less than the last station).
%%                            Branch 3 (e.g. Yes) goes down into an outcome box with the
%%                            given heading and 1 to 3 bullets; branch 4 (e.g. No)
%%                            continues to station N+1. Write \Gate* (with a star) when
%%                            the outcome is a stop or a problem: it adds a warning sign.
%%                            Do not put gates under neighbouring stations.
%%  \SideDocument{A}{B}{Heading}{...}
%%                            Optional, at most two. A file, list or system that the
%%                            process reads (a master list, a register, a template),
%%                            drawn as a document symbol under station A with an arrow
%%                            into station B; heading of 2 to 4 words, 1 to 3 bullets.
%%                            Station A and station B must have no gate; A must not be
%%                            the station right after a gate.
%%  \Loop{F}{T}{text}         Optional, at most two. A return arrow from the bottom of
%%                            station F (normally the last) around the bottom of the
%%                            drawing back into station T, or into the side document under
%%                            station T if there is one. Text: one line (up to 16 words)
%%                            saying when and why. Station F must have no gate and no
%%                            side document under it.
%%  Title block (bottom right; fill from the notes you are given, else leave the line out):
%%  \ChartProject{...}  \ChartDocNo{...}  \ChartRev{...}  \ChartDate{YYYY-MM-DD}
%%  \ChartDrawnBy{...}  \ChartSheet{1 of 1}     (date defaults to today, Rev. to 1)
%%
%%  INSIDE THE BULLET ARGUMENTS ({...} above) you may use only:
%%  \item text               one bullet (the engine adds the list; never \begin{itemize})
%%  \lead{text}              a bold line without a bullet (a heading inside a box)
%%  \plain{text}             a short paragraph without a bullet
%%  \remark{text}            a small italic note (open question, caveat)
%%  \flag{WORD}              a highlighted word, e.g. a status such as \flag{UNMAPPED}
%%  \code{text}              a file, field, button or menu name in typewriter type
%%  \textbf{...} \textit{...}  bold and italic. No nested lists, no \\, no tables.
%%
%%  WRITING RULES
%%  - Facts, not prose: each bullet is one fact in plain words, no full stops at the end.
%%  - Headings: Title Case, 2 to 4 words, verb first for stations (Load Exports, Check
%%    Setup). Branch labels: one word each (Yes/No, Pass/Fail, OK/Problem).
%%  - Keep to the limits above; the engine shrinks the drawing a little if the content
%%    is long, and stamps TOO TALL in red if it cannot fit. If the notes contain more
%%    than fits, keep the facts a reader needs to run the process and drop the rest.
%%  - Characters: write a percent sign as \% and a backslash as \textbackslash. The
%%    signs & _ # $ ^ may be typed as they are. Use plain quotes, no emoji, no symbols
%%    other than arrows (->, written as the word "to" is better).
%%  - Use only the commands listed here. If the process needs something the commands
%%    cannot draw, express it in words inside a bullet rather than inventing a command.
%%  - If the notes are unclear, choose the most sensible reading, draw it, and mention
%%    the choice in one line after the code block. Do not ask questions first.
%%  - Do not invent facts that are not in the notes; do not add company names, people's
%%    names or figures unless they are in the notes.
%%
%% =====================================================================================
%%  PART 2 -- THE ENGINE. Copy unchanged from \documentclass to \begin{document}.
%% =====================================================================================
\documentclass[10pt]{article}
\usepackage[paperwidth=17in, paperheight=11in, margin=0.5in]{geometry}
\usepackage[T1]{fontenc}
\usepackage[utf8]{inputenc}
\usepackage{lmodern}
\usepackage{microtype}
\usepackage{xcolor}
\usepackage{calc}
\usepackage{etoolbox}
\usepackage{pifont}
\usepackage{tikz}
\usetikzlibrary{arrows.meta, calc, shapes.geometric}
\usepackage[hidelinks]{hyperref}
\pagestyle{empty}
\setlength{\parindent}{0pt}
\hyphenpenalty=10000
\exhyphenpenalty=10000
\definecolor{headgreen}{HTML}{548235}
\definecolor{flagpeach}{HTML}{FCE4D6}
\DeclareUnicodeCharacter{2192}{\textrightarrow}
\DeclareUnicodeCharacter{2190}{\textleftarrow}
\DeclareUnicodeCharacter{2713}{\ding{51}}
\DeclareUnicodeCharacter{2714}{\ding{51}}
\DeclareUnicodeCharacter{2717}{\ding{55}}
\DeclareUnicodeCharacter{2718}{\ding{55}}
\DeclareUnicodeCharacter{2022}{\textbullet}
\DeclareUnicodeCharacter{2265}{\ensuremath{\geq}}
\DeclareUnicodeCharacter{2264}{\ensuremath{\leq}}
\DeclareUnicodeCharacter{2260}{\ensuremath{\neq}}
\DeclareUnicodeCharacter{2212}{\textendash}
\DeclareUnicodeCharacter{00D7}{\texttimes}
\DeclareUnicodeCharacter{00B1}{\textpm}
\DeclareUnicodeCharacter{00A0}{~}
\DeclareUnicodeCharacter{2003}{\quad}
\DeclareUnicodeCharacter{20AC}{\texteuro}
\tikzset{
  every node/.style={font=\sffamily\small, inner sep=5pt},
  box/.style={draw, line width=0.5pt, fill=white, align=left},
  flow/.style={-{Stealth[length=9pt, width=6.5pt]}, line width=1.2pt, line join=miter},
  gate/.style={draw, line width=0.8pt, diamond, aspect=2, inner sep=4pt, font=\sffamily\bfseries\small, align=center},
  note/.style={font=\sffamily\footnotesize, inner sep=2pt},
  branch/.style={font=\sffamily\bfseries\small, inner sep=3pt}}
\newcommand{\flag}[1]{{\setlength{\fboxsep}{1pt}\colorbox{flagpeach}{\textbf{#1}}}}
\newcommand{\code}[1]{{\footnotesize\texttt{#1}}}
\makeatletter
% bullets inside boxes: hanging-indent paragraphs (bullet at 0.5em, text at 0.85em, 1.5pt between items)
\newif\iffc@first
\newcommand{\fc@bullet}{\iffc@first\fc@firstfalse\else\par\vskip 1.5pt\fi
  \hangindent=0.85em \hangafter=1 \noindent\hskip 0.85em\llap{\textbullet\hskip 0.35em}\ignorespaces}
\newcommand{\fc@para}[2]{\iffc@first\fc@firstfalse\else\par\vskip 2pt\fi\noindent{#1#2\par}}
\newcommand{\lead}[1]{\fc@para{\bfseries}{#1}}
\newcommand{\plain}[1]{\fc@para{}{#1}}
\newcommand{\remark}[1]{\fc@para{\itshape\fontsize{7.5}{9}\selectfont}{#1}}
% ---- content is recorded by the commands below and drawn at \end{document}
\newcount\fc@n \newcount\fc@ng \newcount\fc@nd \newcount\fc@nl \newcount\fc@i \newcount\fc@j \newcount\fc@k
\newlength{\fc@len}\newlength{\fc@lenb}
\def\fc@title{}\def\fc@subtitle{}\def\fc@sumL{}\def\fc@sumR{}\def\fc@keyhead{}\def\fc@keyitems{}
\def\fc@project{}\def\fc@docno{}\def\fc@rev{1}\def\fc@drawn{}\def\fc@sheet{1 of 1}
\def\fc@date{\the\year-\ifnum\month<10 0\fi\the\month-\ifnum\day<10 0\fi\the\day}
\newcommand{\ChartTitle}[1]{\def\fc@title{#1}\hypersetup{pdftitle={#1}}}
\newcommand{\ChartSubtitle}[1]{\def\fc@subtitle{#1}}
\newcommand{\Summary}[2]{\def\fc@sumL{#1}\def\fc@sumR{#2}}
\newcommand{\KeyBox}[2]{\def\fc@keyhead{#1}\def\fc@keyitems{#2}}
\newcommand{\ChartProject}[1]{\def\fc@project{#1}}
\newcommand{\ChartDocNo}[1]{\def\fc@docno{#1}}
\newcommand{\ChartRev}[1]{\def\fc@rev{#1}}
\newcommand{\ChartDate}[1]{\def\fc@date{#1}}
\newcommand{\ChartDrawnBy}[1]{\def\fc@drawn{#1}}
\newcommand{\ChartSheet}[1]{\def\fc@sheet{#1}}
\newcommand{\Station}[2]{\global\advance\fc@n 1
  \csgdef{fc@sh\the\fc@n}{#1}\csgdef{fc@sb\the\fc@n}{#2}\csgdef{fc@sk\the\fc@n}{0}}
\newcommand{\DocStation}[3]{\global\advance\fc@n 1
  \csgdef{fc@sh\the\fc@n}{#1}\csgdef{fc@sl\the\fc@n}{#2}\csgdef{fc@sb\the\fc@n}{#3}\csgdef{fc@sk\the\fc@n}{1}}
\NewDocumentCommand{\Gate}{s m m m m m m}{\global\advance\fc@ng 1
  \csgdef{fc@gk\the\fc@ng}{#2}\csgdef{fc@gq\the\fc@ng}{#3}\csgdef{fc@gd\the\fc@ng}{#4}\csgdef{fc@go\the\fc@ng}{#5}
  \csgdef{fc@gh\the\fc@ng}{#6}\csgdef{fc@gb\the\fc@ng}{#7}\IfBooleanTF{#1}{\csgdef{fc@gw\the\fc@ng}{1}}{\csgdef{fc@gw\the\fc@ng}{0}}}
\newcommand{\SideDocument}[4]{\global\advance\fc@nd 1
  \csgdef{fc@du\the\fc@nd}{#1}\csgdef{fc@df\the\fc@nd}{#2}\csgdef{fc@dh\the\fc@nd}{#3}\csgdef{fc@db\the\fc@nd}{#4}}
\newcommand{\Loop}[3]{\global\advance\fc@nl 1
  \csgdef{fc@lf\the\fc@nl}{#1}\csgdef{fc@lt\the\fc@nl}{#2}\csgdef{fc@ll\the\fc@nl}{#3}}
\def\doc#1#2{}% only meaningful inside \DocStation; redefined there
% ---- helpers
\newcommand{\fc@items}[1]{\ifblank{#1}{}{\begingroup\let\item\fc@bullet\fc@firsttrue\parindent=0pt \parskip=0pt #1\par\endgroup}}
\newcommand{\fc@stamp}[1]{\listxadd{\fc@stamps}{#1}}\def\fc@stamps{}
\newcommand{\fc@measure}[3]{% #1 width (in, number), #2 font, #3 content -> \fc@len (total height)
  \settototalheight{\fc@len}{\parbox{#1in}{\raggedright #2 #3}}}
\newcommand{\fc@toin}[2]{\pgfmathsetmacro{#1}{(#2)/72.27}}% dimension register -> inches
\newcommand{\fc@max}[2]{\pgfmathsetmacro{#1}{max(#1,#2)}}
\newcommand{\fc@min}[2]{\pgfmathsetmacro{#1}{min(#1,#2)}}
\def\fc@docsym#1#2#3{\draw[line width=0.5pt, fill=white] (#1,#2) -- (#1+#3,#2) -- (#1+#3,#2-0.36)
  .. controls (#1+#3*0.74,#2-0.5) and (#1+#3*0.29,#2-0.26) .. (#1,#2-0.42) -- cycle;}
\def\fc@countdoc#1#2{\global\advance\fc@k 1}
\def\fc@drawdoc#1#2{\fc@docsym{\dx}{\dy}{\dw}
  \node[anchor=north west, align=left, inner sep=0pt, text width=\dwt in] at (\dx+0.07,\dy-0.06)
    {{\footnotesize\texttt{#1}}\ifblank{#2}{}{\\[-1pt]{\fontsize{7.5}{9}\selectfont #2}}};
  \pgfmathsetmacro{\dy}{\dy-0.5}}
\def\fc@addstamp#1{\ifdefempty{\fc@stamptext}{\gdef\fc@stamptext{#1}}{\gappto\fc@stamptext{\\#1}}}
\def\fc@safecatcodes{\catcode`\_=12 \catcode`\&=12 \catcode`\#=12 \catcode`\$=12 \catcode`\^=12 }
% ---- layout: every size in inches on a drawing of W x H units, scaled to 16 x 10 in
\def\fc@layout{%
  \pgfmathsetmacro{\fcW}{12*\fcR}\pgfmathsetmacro{\fcH}{7.5*\fcR}%
  % columns
  \pgfmathsetmacro{\fcw}{min((\fcW-0.5-0.5*(\the\fc@n-1))/\the\fc@n, 2.75)}%
  \pgfmathsetmacro{\fcxz}{0.25+(\fcW-0.5-\the\fc@n*\fcw-0.5*(\the\fc@n-1))/2}%
  \pgfmathsetmacro{\fcwt}{\fcw-10/72.27}%
  \fc@i=0 \loop\ifnum\fc@i<\fc@n \advance\fc@i 1
    \pgfmathsetmacro{\fc@tmp}{\fcxz+(\fc@i-1)*(\fcw+0.5)}\csxdef{fc@xL\the\fc@i}{\fc@tmp}%
    \pgfmathsetmacro{\fc@tmp}{\fc@tmp+\fcw/2}\csxdef{fc@xC\the\fc@i}{\fc@tmp}%
  \repeat
  % header and body heights: the tallest of the row sets both, so every box lines up
  \def\fchh{0.3}\def\fcbh{0.55}%
  \fc@i=0 \loop\ifnum\fc@i<\fc@n \advance\fc@i 1
    \pgfmathsetmacro{\fc@tmp}{\fcw-0.2}%
    \fc@measure{\fc@tmp}{\sffamily\bfseries\normalsize\centering}{\the\fc@i\quad\csuse{fc@sh\the\fc@i}}%
    \fc@toin{\fc@tmp}{\fc@len+8pt}\fc@max{\fchh}{\fc@tmp}%
    \ifnum\csuse{fc@sk\the\fc@i}=1
      \fc@measure{\fcwt}{\sffamily\small}{\csuse{fc@sl\the\fc@i}}\fc@toin{\fc@tmp}{\fc@len}%
      \csxdef{fc@slh\the\fc@i}{\fc@tmp}%
      \fc@k=0 \let\doc\fc@countdoc\setbox0\hbox{\csuse{fc@sb\the\fc@i}}%
      \csxdef{fc@sn\the\fc@i}{\the\fc@k}%
      \pgfmathsetmacro{\fc@tmp}{\fc@tmp+10/72.27+0.14+\the\fc@k*0.5}%
    \else
      \fc@measure{\fcwt}{\sffamily\small}{\fc@items{\csuse{fc@sb\the\fc@i}}}\fc@toin{\fc@tmp}{\fc@len+17pt}%
    \fi
    \fc@max{\fcbh}{\fc@tmp}%
  \repeat
  \pgfmathsetmacro{\fcrowh}{\fchh+\fcbh}%
  % the top zone: title, subtitle, summary columns, key box
  \pgfmathsetmacro{\fctop}{\fcH-0.25}%
  \pgfmathsetmacro{\fcsubY}{\fctop-0.37}%
  \ifdefempty{\fc@subtitle}{\pgfmathsetmacro{\fcsumY}{\fcsubY}}{\pgfmathsetmacro{\fcsumY}{\fcsubY-0.16}}%
  \pgfmathsetmacro{\fczone}{\fcsubY-0.05}%
  \ifdefempty{\fc@sumL}{}{%
    \fc@measure{2.6}{\sffamily\footnotesize}{\fc@items{\fc@sumL}}\fc@toin{\fc@tmp}{\fc@len}%
    \pgfmathsetmacro{\fc@tmp}{\fcsumY-\fc@tmp}\fc@min{\fczone}{\fc@tmp}}%
  \ifdefempty{\fc@sumR}{}{%
    \fc@measure{3.3}{\sffamily\footnotesize}{\fc@items{\fc@sumR}}\fc@toin{\fc@tmp}{\fc@len}%
    \pgfmathsetmacro{\fc@tmp}{\fcsumY-\fc@tmp}\fc@min{\fczone}{\fc@tmp}}%
  \pgfmathsetmacro{\fckw}{\fcW-0.25-6.4}\pgfmathsetmacro{\fckwt}{\fckw-10/72.27}%
  \ifdefempty{\fc@keyhead}{\def\fckh{0}}{%
    \fc@measure{\fckwt}{\sffamily\footnotesize}{\textbf{\small\fc@keyhead}\\[2pt]\fc@items{\fc@keyitems}}%
    \fc@toin{\fckh}{\fc@len+11pt}\pgfmathsetmacro{\fc@tmp}{\fctop-\fckh}\fc@min{\fczone}{\fc@tmp}}%
  \pgfmathsetmacro{\fcrowtop}{\fczone-0.4-\fcshift}%
  \pgfmathsetmacro{\fcrowbot}{\fcrowtop-\fcrowh}%
  % gates: diamond size from the question, outcome box under it
  \def\fclow{\fcrowbot}\def\fcgates{0}\pgfmathsetmacro{\fcow}{\fcw+0.3}\pgfmathsetmacro{\fcowt}{\fcow-10/72.27}%
  \fc@i=0 \loop\ifnum\fc@i<\fc@ng \advance\fc@i 1
    \def\fcgates{1}%
    \settowidth{\fc@len}{\sffamily\bfseries\small\csuse{fc@gq\the\fc@i}}%
    \ifdim\fc@len>1.15in \setlength{\fc@len}{1.15in}\fi
    \fc@toin{\fc@tmp}{\fc@len+1pt}\csxdef{fc@gtw\the\fc@i}{\fc@tmp}%
    \fc@measure{\fc@tmp}{\sffamily\bfseries\small\centering}{\csuse{fc@gq\the\fc@i}}\fc@toin{\fc@tmpb}{\fc@len}%
    \pgfmathsetmacro{\fc@a}{\fc@tmp/2+4/72.27+2*(\fc@tmpb/2+4/72.27)}\csxdef{fc@ga\the\fc@i}{\fc@a}%
    \pgfmathsetmacro{\fc@b}{\fc@a/2}\csxdef{fc@ghb\the\fc@i}{\fc@b}%
    \pgfmathsetmacro{\fc@tmp}{\fcrowbot-0.45-\fc@b}\csxdef{fc@gy\the\fc@i}{\fc@tmp}%
    \pgfmathsetmacro{\fc@tmp}{\fc@tmp-\fc@b-0.45}\csxdef{fc@oy\the\fc@i}{\fc@tmp}%
    \ifnum\csuse{fc@gw\the\fc@i}=1 \pgfmathsetmacro{\fc@hw}{\fcow-0.55}\else\pgfmathsetmacro{\fc@hw}{\fcow-0.18}\fi
    \fc@measure{\fc@hw}{\sffamily\bfseries\normalsize}{\csuse{fc@gh\the\fc@i}}\fc@toin{\fc@hh}{\fc@len}\fc@max{\fc@hh}{0.19}%
    \csxdef{fc@ohh\the\fc@i}{\fc@hh}%
    \fc@measure{\fcowt}{\sffamily\small}{\fc@items{\csuse{fc@gb\the\fc@i}}}\fc@toin{\fc@tmpb}{\fc@len+10pt}%
    \pgfmathsetmacro{\fc@tmpb}{0.12+\fc@hh+\fc@tmpb}\csxdef{fc@oh\the\fc@i}{\fc@tmpb}%
    \pgfmathsetmacro{\fc@tmp}{\fc@tmp-\fc@tmpb}\fc@min{\fclow}{\fc@tmp}%
  \repeat
  \ifnum\fcgates=1 \pgfmathsetmacro{\fcroute}{\fclow-0.25}\fi
  % side documents
  \fc@i=0 \loop\ifnum\fc@i<\fc@nd \advance\fc@i 1
    \pgfmathsetmacro{\fc@tmp}{\fcw-0.16}%
    \fc@measure{\fc@tmp}{\sffamily\small}{\textbf{\csuse{fc@dh\the\fc@i}}\\[2pt]\fc@items{\csuse{fc@db\the\fc@i}}}%
    \fc@toin{\fc@tmp}{\fc@len+22pt}\fc@max{\fc@tmp}{0.7}%
    % does a gate lie between the document's station and the station it feeds? then route below the outcome boxes
    \edef\fc@u{\csuse{fc@du\the\fc@i}}\edef\fc@f{\csuse{fc@df\the\fc@i}}\def\fc@cross{0}%
    {\fc@j=0 \loop\ifnum\fc@j<\fc@ng \advance\fc@j 1
      \edef\fc@kk{\csuse{fc@gk\the\fc@j}}%
      \pgfmathtruncatemacro{\fc@t}{(\fc@kk-\fc@u)*(\fc@kk-\fc@f)<0 ? 1 : 0}%
      \ifnum\fc@t=1 \gdef\fc@cross{1}\fi
    \repeat}%
    \ifnum\fc@cross=1
      \pgfmathsetmacro{\fc@top}{min(\fcrowbot-0.55,\fcroute-0.25+\fc@tmp)}\csxdef{fc@dt\the\fc@i}{\fc@top}%
      \pgfmathsetmacro{\fc@bot}{\fc@top-\fc@tmp}\csxdef{fc@dbot\the\fc@i}{\fc@bot}%
      \csxdef{fc@dr\the\fc@i}{\fcroute}%
    \else
      \pgfmathsetmacro{\fc@top}{\fcrowbot-0.55}\csxdef{fc@dt\the\fc@i}{\fc@top}%
      \pgfmathsetmacro{\fc@bot}{\fc@top-\fc@tmp}\csxdef{fc@dbot\the\fc@i}{\fc@bot}%
      \pgfmathsetmacro{\fc@tmp}{\fc@bot+0.3}\csxdef{fc@dr\the\fc@i}{\fc@tmp}%
    \fi
    \pgfmathsetmacro{\fc@tmp}{\fc@bot-0.1}\fc@min{\fclow}{\fc@tmp}%
  \repeat
  % loops: lanes below everything
  \def\fcbottom{\fclow}%
  \fc@i=0 \loop\ifnum\fc@i<\fc@nl \advance\fc@i 1
    \pgfmathsetmacro{\fc@tmp}{\fclow-0.45-(\fc@i-1)*0.4}\csxdef{fc@ly\the\fc@i}{\fc@tmp}\def\fcbottom{\fc@tmp}%
  \repeat
  \pgfmathsetmacro{\fcslack}{\fcbottom-0.65}%
}
% ---- draw
\def\fc@draw{%
  \ifnum\fc@n<2 \fc@stamp{Add at least two \string\Station\ blocks}\fi
  \ifnum\fc@n>6 \fc@stamp{At most six stations: merge or drop some}\fi
  \ifdefempty{\fc@title}{\fc@stamp{MISSING: \string\ChartTitle}\def\fc@title{\textcolor{red}{MISSING TITLE}}}{}%
  \ifnum\fc@n<1 \global\advance\fc@n 1 \csgdef{fc@sh1}{Station}\csgdef{fc@sb1}{\item Add stations}\csgdef{fc@sk1}{0}\fi
  \def\fcR{1}\def\fcshift{0}\fc@layout
  \ifdim\fcslack pt<0pt
    \pgfmathsetmacro{\fcR}{min(1.45,(7.5-\fcslack)/7.5*1.02)}\fc@layout
    \ifdim\fcslack pt<0pt \pgfmathsetmacro{\fcR}{min(1.45,\fcR*(7.5*\fcR-\fcslack)/(7.5*\fcR)*1.02)}\fc@layout\fi
    \ifdim\fcslack pt<0pt \fc@stamp{TOO TALL: shorten the bullets or drop a station}\fi
  \fi
  \ifdim\fcslack pt>0pt \pgfmathsetmacro{\fcshift}{min(\fcslack/2,0.9)}\fc@layout\fi
  \settowidth{\fc@len}{\sffamily\bfseries\footnotesize\fc@title}\fc@toin{\fctw}{\fc@len+12pt}\fc@max{\fctw}{2.6}%
  \pgfmathsetmacro{\fcscale}{1.333333/\fcR}%
  \noindent\scalebox{\fcscale}{%
  \begin{tikzpicture}[x=1in, y=1in]
  \useasboundingbox (0,0) rectangle (\fcW,\fcH);
  \draw[line width=0.8pt] (0,0) rectangle (\fcW,\fcH);
  % title zone
  \node[anchor=north west, inner sep=0pt, font=\sffamily\bfseries\fontsize{14}{16}\selectfont] at (0.25,\fctop) {\fc@title};
  \ifdefempty{\fc@subtitle}{}{\node[anchor=north west, inner sep=0pt, font=\sffamily\footnotesize\itshape] at (0.25,\fcsubY) {\fc@subtitle};}
  \ifdefempty{\fc@sumL}{}{\node[anchor=north west, align=left, inner sep=0pt, text width=2.6in, font=\sffamily\footnotesize] at (0.25,\fcsumY) {\fc@items{\fc@sumL}};}
  \ifdefempty{\fc@sumR}{}{\node[anchor=north west, align=left, inner sep=0pt, text width=3.3in, font=\sffamily\footnotesize] at (2.95,\fcsumY) {\fc@items{\fc@sumR}};}
  \ifdefempty{\fc@keyhead}{}{\node[box, anchor=north west, text width=\fckwt in, font=\sffamily\footnotesize] at (6.4,\fctop) {\textbf{\small\fc@keyhead}\\[2pt]\fc@items{\fc@keyitems}};}
  % stations
  \foreach \i in {1,...,\the\fc@n} {
    \edef\xL{\csuse{fc@xL\i}}\pgfmathsetmacro{\xR}{\xL+\fcw}\pgfmathsetmacro{\yb}{\fcrowtop-\fchh}\pgfmathsetmacro{\yc}{\fcrowtop-\fchh/2}
    \pgfmathsetmacro{\ybb}{\yb-\fcbh}
    \draw[line width=0.5pt, fill=headgreen] (\xL,\fcrowtop) rectangle (\xR,\yb);
    \pgfmathsetmacro{\hw}{\fcw-0.2}
    \node[text=white, font=\sffamily\bfseries\normalsize, inner sep=0pt, text width=\hw in, align=center] at (\xL+\fcw/2,\yc) {\i\quad\csuse{fc@sh\i}};
    \draw[box] (\xL,\yb) rectangle (\xR,\ybb);
    \ifnum\csuse{fc@sk\i}=1
      \node[anchor=north west, inner sep=5pt, text width=\fcwt in, align=left] at (\xL,\yb) {\csuse{fc@sl\i}};
      \pgfmathsetmacro{\dw}{\fcw-0.34}\pgfmathsetmacro{\dwt}{\dw-0.14}\pgfmathsetmacro{\dx}{\xL+0.17}
      \pgfmathsetmacro{\dy}{\yb-10/72.27-\csuse{fc@slh\i}-0.1}
      \let\doc\fc@drawdoc
      \csuse{fc@sb\i}
    \else
      \node[anchor=north west, inner sep=5pt, text width=\fcwt in, align=left] at (\xL,\yb) {\fc@items{\csuse{fc@sb\i}}};
    \fi
    \ifnum\i<\fc@n \draw[flow] (\xR,\yc) -- (\xR+0.5,\yc); \fi
  }
  % gates
  \ifnum\fc@ng>0 \foreach \g in {1,...,\the\fc@ng} {
    \edef\k{\csuse{fc@gk\g}}
    \ifnum\k<1 \fc@stamp{Gate \g: station number \k\ does not exist}\else\ifnum\k>\fc@n \fc@stamp{Gate \g: station number \k\ does not exist}\else
    \ifnum\k=\fc@n \fc@stamp{Gate under the last station: move it to an earlier station or add a final station}\fi
    \edef\gx{\csuse{fc@xC\k}}\edef\gy{\csuse{fc@gy\g}}\edef\ga{\csuse{fc@ga\g}}\edef\gb{\csuse{fc@ghb\g}}
    \edef\oy{\csuse{fc@oy\g}}\edef\oh{\csuse{fc@oh\g}}\edef\ohh{\csuse{fc@ohh\g}}\pgfmathsetmacro{\oit}{\oy-0.1-\ohh}
    \draw[flow] (\gx,\fcrowbot) -- (\gx,\gy+\gb);
    \edef\gtw{\csuse{fc@gtw\g}}
    \node[gate, text width=\gtw in] at (\gx,\gy) {\csuse{fc@gq\g}};
    \draw[flow] (\gx,\gy-\gb) -- (\gx,\oy) node[midway, right, branch] {\csuse{fc@gd\g}};
    \draw[box] (\gx-\fcow/2,\oy) rectangle (\gx+\fcow/2,\oy-\oh);
    \ifnum\csuse{fc@gw\g}=1
      \draw[line width=0.8pt, fill=yellow!85, rounded corners=1pt] (\gx-\fcow/2+0.08,\oy-0.3) -- ++(0.3,0) -- ++(-0.15,0.26) -- cycle;
      \node[font=\sffamily\bfseries\small, inner sep=0pt] at (\gx-\fcow/2+0.23,\oy-0.2) {!};
      \pgfmathsetmacro{\ohw}{\fcow-0.55}
      \node[anchor=north west, font=\sffamily\bfseries\normalsize, inner sep=0pt, text width=\ohw in, align=left] at (\gx-\fcow/2+0.45,\oy-0.09) {\csuse{fc@gh\g}};
    \else
      \pgfmathsetmacro{\ohw}{\fcow-0.18}
      \node[anchor=north west, font=\sffamily\bfseries\normalsize, inner sep=0pt, text width=\ohw in, align=left] at (\gx-\fcow/2+0.08,\oy-0.09) {\csuse{fc@gh\g}};
    \fi
    \node[anchor=north west, inner sep=5pt, text width=\fcowt in, align=left] at (\gx-\fcow/2,\oit) {\fc@items{\csuse{fc@gb\g}}};
    \ifnum\k<\fc@n
      \pgfmathtruncatemacro{\kk}{\k+1}\edef\nx{\csuse{fc@xC\kk}}
      \draw[line width=1.2pt] (\gx+\ga,\gy) -- (\nx,\gy) node[pos=0.5, above, branch] {\csuse{fc@go\g}};
      \ifcsdef{fc@fed\kk}{\fill (\nx,\gy) circle (2.6pt);}{\draw[flow] (\nx,\gy) -- (\nx,\fcrowbot);}
    \fi
    \fi\fi
  }\fi
  % side documents
  \ifnum\fc@nd>0 \foreach \d in {1,...,\the\fc@nd} {
    \edef\u{\csuse{fc@du\d}}\edef\f{\csuse{fc@df\d}}
    \ifnum\u<1 \fc@stamp{Side document \d: station number \u\ does not exist}\else\ifnum\u>\fc@n \fc@stamp{Side document \d: station number \u\ does not exist}\else
    \ifnum\f<1 \fc@stamp{Side document \d: station number \f\ does not exist}\else\ifnum\f>\fc@n \fc@stamp{Side document \d: station number \f\ does not exist}\else
    \edef\xL{\csuse{fc@xL\u}}\pgfmathsetmacro{\xR}{\xL+\fcw}\edef\xC{\csuse{fc@xC\u}}
    \edef\dt{\csuse{fc@dt\d}}\edef\db{\csuse{fc@dbot\d}}\edef\dr{\csuse{fc@dr\d}}\pgfmathsetmacro{\dwt}{\fcw-0.16}
    \pgfmathsetmacro{\dfoot}{\db-0.039}\csxdef{fc@under\u}{\dfoot}
    \draw[line width=0.5pt, fill=white] (\xL,\dt) -- (\xR,\dt) -- (\xR,\db)
      .. controls (\xL+0.65*\fcw,\db-0.16) and (\xL+0.35*\fcw,\db+0.08) .. (\xL,\db-0.07) -- cycle;
    \node[anchor=north west, align=left, inner sep=0pt, text width=\dwt in] at (\xL+0.08,\dt-0.08) {\textbf{\csuse{fc@dh\d}}\\[2pt]\fc@items{\csuse{fc@db\d}}};
    \edef\fx{\csuse{fc@xC\f}}
    \ifnum\f>\u \draw[flow] (\xR,\dr) -| (\fx,\fcrowbot); \fi
    \ifnum\f<\u \draw[flow] (\xL,\dr) -| (\fx,\fcrowbot); \fi
    \ifnum\f=\u \draw[flow] (\xC,\dt) -- (\xC,\fcrowbot); \fi
    \fi\fi\fi\fi
  }\fi
  % loops
  \ifnum\fc@nl>0 \foreach \l in {1,...,\the\fc@nl} {
    \edef\f{\csuse{fc@lf\l}}\edef\t{\csuse{fc@lt\l}}\edef\ly{\csuse{fc@ly\l}}
    \ifnum\f<1 \fc@stamp{Loop \l: station number \f\ does not exist}\else\ifnum\f>\fc@n \fc@stamp{Loop \l: station number \f\ does not exist}\else
    \ifnum\t<1 \fc@stamp{Loop \l: station number \t\ does not exist}\else\ifnum\t>\fc@n \fc@stamp{Loop \l: station number \t\ does not exist}\else
    \edef\fx{\csuse{fc@xC\f}}\edef\tx{\csuse{fc@xC\t}}
    \ifcsdef{fc@busy\f}{\pgfmathsetmacro{\fx}{\fx+\fcw/4}}{}
    \ifcsdef{fc@under\t}{\edef\ey{\csuse{fc@under\t}}}{\edef\ey{\fcrowbot}\ifcsdef{fc@busy\t}{\pgfmathsetmacro{\tx}{\tx-\fcw/4}}{}}
    \draw[flow] (\fx,\fcrowbot) |- (\tx,\ly) -- (\tx,\ey);
    \node[anchor=south, note] at ({(\fx+\tx)/2},\ly+0.03) {\csuse{fc@ll\l}};
    \fi\fi\fi\fi
  }\fi
  % title block, one row in the bottom right corner
  \pgfmathsetmacro{\tbx}{\fcW-5.05-\fctw}
  \pgfmathsetmacro{\tba}{\fcW-5.05}\pgfmathsetmacro{\tbb}{\tba+1.4}\pgfmathsetmacro{\tbc}{\tbb+0.9}
  \pgfmathsetmacro{\tbd}{\tbc+0.45}\pgfmathsetmacro{\tbe}{\tbd+0.85}\pgfmathsetmacro{\tbf}{\tbe+0.9}
  \draw[line width=0.5pt] (\tbx,0) -- (\tbx,0.42) -- (\fcW,0.42);
  \foreach \x in {\tba,\tbb,\tbc,\tbd,\tbe,\tbf} \draw[line width=0.5pt] (\x,0) -- (\x,0.42);
  \foreach \x/\lab/\val in {\tbx/Title/\fc@title, \tba/Project/\fc@project, \tbb/Doc no./\fc@docno, \tbc/Rev./\fc@rev,
                            \tbd/Date/\fc@date, \tbe/Drawn/\fc@drawn, \tbf/Sheet/\fc@sheet} {
    \node[anchor=north west, inner sep=0pt, font=\sffamily\fontsize{7}{8}\selectfont] at (\x+0.07,0.36) {\lab};
    \node[anchor=south west, inner sep=0pt, font=\sffamily\bfseries\footnotesize] at (\x+0.07,0.07) {\val};
  }
  % stamps: red lines bottom left when something could not be drawn as asked
  \gdef\fc@stamptext{}\forlistloop{\fc@addstamp}{\fc@stamps}
  \ifdefempty{\fc@stamptext}{}{\node[anchor=south west, inner sep=0pt, text=red, font=\sffamily\bfseries\footnotesize, text width=3.9in, align=left] at (0.3,0.1) {\fc@stamptext};}
  \end{tikzpicture}}%
}
\def\fc@checks{% layout rules a weak layout cannot draw; each becomes a red stamp
  \fc@i=0 \loop\ifnum\fc@i<\fc@ng \advance\fc@i 1
    \edef\k{\csuse{fc@gk\the\fc@i}}\csgdef{fc@gated\k}{1}\pgfmathtruncatemacro{\fc@tmp}{\k+1}\csgdef{fc@busy\fc@tmp}{1}%
    {\fc@j=\fc@i \loop\ifnum\fc@j<\fc@ng \advance\fc@j 1
      \edef\kk{\csuse{fc@gk\the\fc@j}}\pgfmathtruncatemacro{\fc@tmp}{abs(\k-\kk)}%
      \ifnum\fc@tmp<2 \fc@stamp{Gates under stations \k\ and \kk: keep at least one station between gates}\fi
    \repeat}
  \repeat
  \fc@i=0 \loop\ifnum\fc@i<\fc@nd \advance\fc@i 1
    \edef\u{\csuse{fc@du\the\fc@i}}\edef\f{\csuse{fc@df\the\fc@i}}%
    \csgdef{fc@fed\f}{1}\csgdef{fc@hasdoc\u}{1}\csgdef{fc@busy\f}{1}%
    \ifcsdef{fc@gated\u}{\fc@stamp{Side document under station \u, which has a gate: put it under another station}}{}%
    \ifcsdef{fc@gated\f}{\fc@stamp{Side document feeds station \f, which has a gate: feed another station}}{}%
    \pgfmathtruncatemacro{\fc@tmp}{\u-1}\ifcsdef{fc@gated\fc@tmp}{\fc@stamp{Side document under station \u, right after a gate: put it under another station}}{}%
  \repeat
  \fc@i=0 \loop\ifnum\fc@i<\fc@nl \advance\fc@i 1
    \edef\f{\csuse{fc@lf\the\fc@i}}\edef\t{\csuse{fc@lt\the\fc@i}}%
    \ifcsdef{fc@gated\f}{\fc@stamp{Loop from station \f, which has a gate: start the loop at another station}}{}%
    \ifcsdef{fc@hasdoc\f}{\fc@stamp{Loop from station \f, which has a side document under it: start it elsewhere}}{}%
    \ifcsdef{fc@gated\t}{\fc@stamp{Loop into station \t, which has a gate: end the loop at another station}}{}%
  \repeat
}
\AtEndDocument{\fc@checks\fc@draw}
\makeatother
\begin{document}
\makeatletter\fc@safecatcodes\makeatother
%% =====================================================================================
%%  PART 3 -- CONTENT. Replace every line below down to \end{document}.
%%  (As shipped, the placeholders draw a legend of the building blocks.)
%% =====================================================================================
\ChartTitle{Chart Title -- Document Type}
\ChartSubtitle{One italic line on what this chart shows and who it is for.}
\Summary{\item Key fact one, up to ten words \item Key fact two \item Key fact three}
        {\item Key fact four \item Key fact five}
\KeyBox{Key box heading}{\item A precondition or rule, up to fourteen words
  \item A second precondition \item A third, with a \code{file or button name}}
\DocStation{Document Station}{A lead line}{\doc{Input file name.xlsx}{required}
  \doc{Second file name.csv}{optional}}
\Station{Station Heading}{\item First fact, up to twelve words \item Second fact
  \item Third fact with a \flag{STATUS} word \remark{A small italic remark or open question.}}
\Station{Station With Gate}{\item A station that ends in a question \item Its gate hangs below
  \item The branch that continues goes to the next station}
\Station{Station Heading}{\lead{A bold lead line} \item First fact \item Second fact
  \plain{A short plain paragraph without a bullet, for a result or a rule.}}
\Station{Last Station}{\item First fact \item Second fact \item The loop starts here}
\Gate*{3}{Question?}{Yes}{No}{Outcome Heading}{\item What happens on this branch \item A second fact}
\SideDocument{1}{4}{Side Document}{\item A list or file the process reads \item Which station uses it}
\Loop{5}{1}{Loop text: when and why the process returns to the start}
\ChartProject{Project name}
\ChartDocNo{DOC-000}
\ChartRev{1}
\ChartDrawnBy{Your name}
\ChartSheet{1 of 1}
\end{document}
